% [VERIFY] methodology traces to: concepts.md; solution/algorithm.md (Eqn1,Eqn2,pseudocode); solution/architecture.md; solution/constraints.md; solution/heuristics.md; Table 1; C02,C04,C05
\section{Deep Residual Learning}
\label{sec:method}

\subsection{Residual learning}
\label{sec:method:residual}
Let $\Hmap(\xin)$ denote an underlying mapping to be fit by a few stacked layers, with
$\xin$ the input to the first of these layers. If one hypothesizes that multiple
nonlinear layers can asymptotically approximate a complicated function, then it is
equivalent to hypothesize that they can approximate the residual function
$\Fmap(\xin) := \Hmap(\xin) - \xin$ (assuming input and output have the same
dimension). Rather than expecting the stacked layers to approximate $\Hmap(\xin)$, we
explicitly let them approximate $\Fmap(\xin)$, and recover the original mapping as
$\Fmap(\xin) + \xin$. Although both forms should be able to approximate the desired
function asymptotically, the ease of learning differs. This reformulation is motivated
by the counter-intuitive degradation phenomenon discussed in \cref{sec:intro}: if added
layers can be constructed as identity mappings, a deeper model should have training
error no greater than its shallower counterpart, yet solvers struggle to find such a
solution. With residual learning, if an identity mapping is optimal, the solver may
simply drive the weights of the nonlinear layers toward zero to approach identity,
rather than fit an identity by a stack of nonlinear layers. In practice the optimal
functions are unlikely to be exactly identity, but the reformulation provides a
better-conditioned starting point: if optimal functions are closer to identity than to
zero, it is easier for the solver to find perturbations with reference to an identity
mapping than to learn the function afresh.

\subsection{Identity mapping by shortcuts}
\label{sec:method:shortcut}
We adopt residual learning to every few stacked layers. A building block is defined as
\begin{equation}
  \yout = \Fmap(\xin, \{W_i\}) + \xin, \label{eq:block}
\end{equation}
where $\xin$ and $\yout$ are the input and output vectors of the layers considered and
$\Fmap(\xin, \{W_i\})$ is the residual mapping to be learned. The operation
$\Fmap + \xin$ is performed by a shortcut connection and an element-wise addition,
followed by a second nonlinearity $\sigma(\yout)$. For a two-layer block,
$\Fmap = W_2\,\sigma(W_1 \xin)$ with $\sigma$ the rectified linear unit and biases
omitted; a three-layer block applies $\Fmap = W_3\,\sigma(W_2\,\sigma(W_1 \xin))$. The
shortcut connection in \cref{eq:block} introduces neither extra parameters nor
computational complexity beyond a negligible element-wise addition. This is not only
attractive in practice but also central to the comparisons: the plain and residual
networks then have the same number of parameters, depth, width, and computational cost,
so a comparison between them is controlled.

\Cref{eq:block} requires $\Fmap(\xin)$ and $\xin$ to share dimensions. When they differ,
for instance where the number of channels changes, a linear projection $\Wproj$ is
applied on the shortcut to match dimensions,
\begin{equation}
  \yout = \Fmap(\xin, \{W_i\}) + \Wproj\,\xin. \label{eq:projblock}
\end{equation}
A square $\Wproj$ may also be used in \cref{eq:block}, but experiments
(\cref{sec:exp:shortcut}) show that identity shortcuts are sufficient for addressing the
\degprob and are economical, so the projection in \cref{eq:projblock} is reserved for
matching dimensions. The residual function $\Fmap$ is flexible, but a single-layer
$\Fmap$ is similar to a linear layer and offers no observed advantage, so we use
$\Fmap$ with two or more layers.
% METHOD-MORE

\subsection{Network architectures}
\label{sec:method:arch}
We instantiate residual learning in a stage-structured convolutional network whose plain
baseline follows VGG-style design rules: layers that produce feature maps of the same
size keep the same number of filters, and when the spatial size is halved the number of
filters is doubled so that the time complexity per layer is preserved. A single stem
applies a $7\times7$ convolution with 64 filters at stride 2, followed by batch
normalization, a rectified linear unit, and a $3\times3$ max-pool at stride 2, producing
a $56\times56\times64$ feature map shared by every depth. Four residual stages
(\texttt{conv2\_x} through \texttt{conv5\_x}) then stack building blocks; the first block
of \texttt{conv3\_x}, \texttt{conv4\_x}, and \texttt{conv5\_x} performs down-sampling
with a stride-2 convolution rather than pooling, and the shortcut at those stage
boundaries down-samples to match. A global average pool, a 1000-way fully-connected
layer, and a softmax form the head. The residual network is obtained from the plain
network by inserting shortcuts (\cref{eq:block}) between every pair (basic block) or
triple (bottleneck block) of layers, leaving the plain skeleton otherwise unchanged.

For the basic block used at 18 and 34 layers, $\Fmap$ is two $3\times3$ convolutions. To
reach greater depth economically we use a \emph{bottleneck} block: a $1\times1$
convolution reduces the channel dimension, a $3\times3$ convolution operates on the
reduced bottleneck, and a second $1\times1$ convolution restores the high-dimensional
output, with the output width four times the bottleneck width. A bottleneck block carries
the same per-block time complexity as the basic two-layer block while spanning three
weighted layers, which is what makes 50/101/152-layer networks tractable. Identity
shortcuts are especially important in bottleneck designs, because replacing the identity
on the high-dimensional ends with a projection would roughly double both the time
complexity and the model size of the block. \Cref{tab:arch} summarizes the per-depth
stage layouts and floating-point complexity; the 34-layer model has the same complexity
as its plain counterpart, and the 152-layer model, while eight times deeper than VGG-19,
still has lower complexity.

% Source: evidence/tables/table1_imagenet_architectures.md (Table 1). Supports C03, C05.
% EXACT numbers. FLOPs in units of 1e9.
\begin{table}[t]
\centering
\caption{Architectures for ImageNet. Building blocks are shown in brackets with the
number of stacked blocks. Down-sampling is performed by \texttt{conv3\_1},
\texttt{conv4\_1}, and \texttt{conv5\_1} with a stride of 2. Source: Table~1 of the
underlying artifact; supports the depth-scaling and bottleneck claims.}
\label{tab:arch}
\small
\setlength{\tabcolsep}{4pt}
\resizebox{\textwidth}{!}{%
\begin{tabular}{l l c c c}
\toprule
stage & output & 18 / 34-layer (basic) & 50 / 101 / 152-layer (bottleneck) & blocks (18/34/50/101/152) \\
\midrule
\texttt{conv1}   & $112\times112$ & \multicolumn{2}{c}{$7\times7$, 64, stride 2} & --- \\
\texttt{conv2\_x} & $56\times56$  & $[3{\times}3,\,64]\times2$ & $[1{\times}1,64;\,3{\times}3,64;\,1{\times}1,256]$ & 2 / 3 / 3 / 3 / 3 \\
\texttt{conv3\_x} & $28\times28$  & $[3{\times}3,\,128]\times2$ & $[1{\times}1,128;\,3{\times}3,128;\,1{\times}1,512]$ & 2 / 4 / 4 / 4 / 8 \\
\texttt{conv4\_x} & $14\times14$  & $[3{\times}3,\,256]\times2$ & $[1{\times}1,256;\,3{\times}3,256;\,1{\times}1,1024]$ & 2 / 6 / 6 / 23 / 36 \\
\texttt{conv5\_x} & $7\times7$    & $[3{\times}3,\,512]\times2$ & $[1{\times}1,512;\,3{\times}3,512;\,1{\times}1,2048]$ & 2 / 3 / 3 / 3 / 3 \\
                 & $1\times1$    & \multicolumn{2}{c}{average pool, 1000-d fc, softmax} & --- \\
\midrule
FLOPs ($\times10^{9}$) & & 1.8 (18), 3.6 (34) & 3.8 (50), 7.6 (101), 11.3 (152) & VGG-19: 19.6 \\
\bottomrule
\end{tabular}%
}
\end{table}


For dimension-changing shortcuts we consider three options: \optA{} uses identity
mapping with zero-padding for the increased channels and is entirely parameter-free;
\optB{} uses projection shortcuts (\cref{eq:projblock}) only where dimensions change and
identity elsewhere; and \optC{} uses projections on all shortcuts. \Cref{sec:exp:shortcut}
quantifies their small differences. Within a stage, where dimensions do not change,
identity shortcuts are always used. On CIFAR-10 we use a separate, deliberately minimal
family to study depth in isolation: a first $3\times3$ convolution with 16 filters,
followed by $2n$ layers at each of three feature-map sizes $\{32,16,8\}$ with widths
$\{16,32,64\}$, giving $6n{+}2$ weighted layers, identity (\optA) shortcuts everywhere,
and a global-average-pool/softmax head. Because the CIFAR shortcuts are all identity,
the residual nets have \emph{exactly} the same depth, width, and parameter count as their
plain counterparts.

\subsection{Implementation and training}
\label{sec:method:impl}
We follow a single recipe across plain and residual variants so that comparisons are
controlled. Batch normalization is applied right after each convolution and before
activation; weights are initialized with the variance-preserving scheme for
rectifiers; and no dropout is used. Networks are trained from scratch with stochastic
gradient descent, momentum $0.9$, and weight decay $10^{-4}$. The complete forward pass
of a single residual block is given in \cref{alg:block}; the full network applies it
stage by stage between the stem and the head. Design choices in this section are tied to
the experimental verifications in \cref{sec:exp}: the controlled plain-versus-residual
comparison (\cref{sec:exp:imagenet}), the shortcut-option ablation
(\cref{sec:exp:shortcut}), the bottleneck depth scan (\cref{sec:exp:depth}), and the
CIFAR-10 and detection-transfer studies (\cref{sec:exp:cifar,sec:exp:detection}).

\begin{algorithm}[t]
\caption{Forward pass of one residual building block}
\label{alg:block}
\begin{algorithmic}[1]
\Function{ResidualBlock}{$\xin$, $\Fmap$, option, downsample}
  \State $\mathrm{out} \gets \Fmap(\xin)$ \Comment{conv\,$\rightarrow$\,BN\,$\rightarrow$\,ReLU\,$\rightarrow$\,conv\,$\rightarrow$\,BN (or $1{\times}1$\,$\rightarrow$\,$3{\times}3$\,$\rightarrow$\,$1{\times}1$)}
  \If{$\dim(\mathrm{out}) = \dim(\xin)$ \textbf{and not} downsample}
    \State $\mathrm{identity} \gets \xin$ \Comment{parameter-free identity shortcut}
  \ElsIf{option $=$ A}
    \State $\mathrm{identity} \gets \textsc{ZeroPadChannels}(\textsc{MaybeStride2}(\xin))$
  \Else
    \State $\mathrm{identity} \gets \textsc{Projection}(\xin)$ \Comment{$1{\times}1$ conv, optional stride 2}
  \EndIf
  \State $\mathrm{out} \gets \mathrm{out} + \mathrm{identity}$ \Comment{element-wise addition}
  \State \Return $\textsc{ReLU}(\mathrm{out})$
\EndFunction
\end{algorithmic}
\end{algorithm}

